% ECONOMICS_PROBLEM: {"id": "C73-40", "title": "Nonlinear stability of periodic MFG solutions", "jel_code": "C73", "source_id": "OP-0211"}
% FIRST_POSTED: 2026-10-10
% ATTEMPT_STATUS: solved
% ENTRY_KIND: research_attempt
% !TEX program = pdflatex
\documentclass[11pt,a4paper]{article}
\usepackage[T1]{fontenc}
\usepackage{lmodern}
\usepackage[left=20mm,right=20mm,top=18mm,bottom=18mm]{geometry}
\usepackage{amsmath,amssymb,amsthm,mathtools,microtype}
\usepackage[colorlinks=true,linkcolor=blue,citecolor=blue,urlcolor=blue,hypertexnames=false]{hyperref}
\allowdisplaybreaks
\setlength{\parindent}{0pt}
\setlength{\parskip}{3pt}
\setlength{\emergencystretch}{2em}
\newtheorem{theorem}{Theorem}
\newtheorem{lemma}[theorem]{Lemma}
\newtheorem{corollary}[theorem]{Corollary}
\title{Proof: Selected boundary stability for periodic mean-field games}
\author{GPT Codex}
\date{10 October 2026}
\begin{document}
\maketitle
\begin{abstract}
For every dimension, smooth local coupling and classical genuinely
periodic quadratic mean-field-game equilibrium, we characterize a
specified weaker nonlinear stability property: all small smooth
boundary perturbations admit a nearby selected equilibrium with an
error linear in the boundary error, with constants depending on the
horizon. The criterion is both necessary and sufficient, including a
negative conclusion at every conjugate horizon. It has finite matrix
certificates constructed from two scalar parabolic problems, and a
convergent equilibrium-selection rule. An explicit entire periodic
family has only a discrete set of conjugate horizons. All arguments
are given below.
\end{abstract}

\textbf{Scope.} This solves the expressly permitted weaker-stability
alternative in C73-40 for its entire original model class. The property
is defined before its criterion. No restriction to a special coupling,
dimension, periodic family or a preselected set of regular horizons is
made in the general theorem. The original horizon-uniform orbital
shadowing of every admissible equilibrium is a stronger question and
is not claimed here.

The complete original Proof is preserved at the immutable
\href{https://github.com/jbaelaw/periodic-mfg-selected-boundary-stability/blob/795f26f0dc1ec0c9d5838156779cffa692ede0cf/paper/Proof.tex}{Proof TeX}
and \href{https://github.com/jbaelaw/periodic-mfg-selected-boundary-stability/blob/795f26f0dc1ec0c9d5838156779cffa692ede0cf/paper/Proof.pdf}{Proof PDF}.

\section{The original boundary-value setting}
Let $\nu>0$, let $f:\mathbb T^d\times(0,\infty)\to\mathbb R$ be
smooth, and let a positive, mass-one classical reference satisfy
\begin{equation}\tag{MFG}\label{MFG}
 \begin{split}
 \lambda_*-v_{*,t}-\nu\Delta v_*+\tfrac12|Dv_*|^2&=f(x,m_*),\\
 m_{*,t}-\nu\Delta m_*-\operatorname{div}(m_*Dv_*)&=0,\\
 (v_*,m_*)(t+P)&=(v_*,m_*)(t),\qquad \int m_*(t)=1.
 \end{split}
\end{equation}
The pair $(Dv_*,m_*)$ is nonconstant in time. The actual reference
value $v_*(t,x)-\lambda_*t$ may drift. Its orbit is
$\mathcal O_* =\{(Dv_*(s),m_*(s)):s\in\mathbb R/P\mathbb Z\}$,
and its distance is the infimum over $s$ of the sum of the gradient
$L^\infty$ error and density $L^1$ error.

Perturbed equilibria solve
\[
 -u_t-\nu\Delta u+\tfrac12|Du|^2=f(x,m),\qquad
 m_t-\nu\Delta m-\operatorname{div}(mDu)=0,
 \qquad m(0)=m_0,\quad u(T)=h.
\]
Initial population and terminal value are the independent boundary
data. No initial-value MFG semigroup is presumed.

Cardaliaguet--Delarue discuss the periodic-solution stability question
in \cite[Section 3.3, paragraph ``Periodic solutions,'' p.~3681]{CD}.
Genuine periodic examples already appear in
\cite[Section 6, Theorem 6.1 and Remark 6.4]{CN}; the construction
below is not a claim of first existence. The classification proved
here concerns the specified nonlinear selection property, rather
than imaginary eigenvalues or an unverified uniform inverse.

% Research component, not a published or independently approved final paper.
\section{The nonlinear property being classified}
Write $Q_T=[0,T]\times\mathbb T^d$, with the torus of volume one, and
$\Pi z=z-\int_{\mathbb T^d}z$. Fix $0<\alpha<1$.
For any reference in the original problem, any $T>0$, and any phase
$\theta$, put
\[
 \bar m(t)=m_*(t+\theta),\quad b(t)=Dv_*(t+\theta),\quad
 a(t,x)=f_m(x,\bar m(t,x)),\quad
 u_{\rm r}(t,x)=v_*(t+\theta,x)-\lambda_*t.
\]
Both $\bar m$ and $b$ are known reference fields, not unknown solutions
of a perturbed boundary problem. We use
\[
 X_T=\{(w,\rho)\in(C^{1+\alpha/2,2+\alpha}(Q_T))^2:
                  \textstyle\int w(t)=\int\rho(t)=0\},\qquad
 D=(C^{2+\alpha}_0(\mathbb T^d))^2,
\]
with their usual sum norms; a subscript zero means spatial mean zero.
The boundary perturbation is
$d=(d_0,d_T)=(m_0-\bar m(0),\Pi(h-u_{\rm r}(T)))$.
Every sufficiently small $d\in D$ gives a positive normalized initial
density. Terminal constants are free and will be restored explicitly.

\textbf{Selected local Lipschitz boundary stability at $(T,\theta)$.}
There are $r,\delta,C>0$ and a rule $S:\{\|d\|_D<\delta\}\to X_T$
such that $S(0)=0$, and for every such $d$, $S(d)=(w,\rho)$ produces
a classical MFG equilibrium with the stated initial density and terminal
value, with
\begin{equation}\tag{W1}\label{W1}
 \|(w,\rho)\|_{X_T}\leq C\|d\|_D<r,
 \qquad \bar m+\rho\geq\tfrac12\min_{Q_T}\bar m.
\end{equation}
No continuity or differentiability of $S$ is assumed in this definition.
The theorem constructs a Lipschitz rule when the property
holds. This is a specified weaker alternative in C73-40: boundary errors
are measured more strongly, a nearby equilibrium is selected, and the
constants may depend on horizon and phase. In particular,
\begin{equation}\tag{W2}\label{W2}
 \sup_{0\leq t\leq T}d_*((Du(t),m(t)),\mathcal O_*)
 \leq\|(w,\rho)\|_{X_T}\leq C\|d\|_D.
\end{equation}
The comparison phase in this estimate is $t+\theta$ modulo $P$.
The neighborhood can always be chosen with density lower bound
$\min m_*/2$ and an admissibility upper bound one larger than the
reference norm. If the original specified caps have strict slack, it
can also be chosen within those caps. We do not assert an open
perturbation result inside an arbitrarily saturated cap.

For a set $H\subset(0,\infty)\times(\mathbb R/P\mathbb Z)$, the
reference has this property on $H$ when it has it at every pair in $H$.
These outer quantifiers are part of the property; the criterion below
classifies every pair, including the pairs where the property fails.

\section{Scalar parabolic facts and an explicit compactness bound}
We record the analytic facts used in the classification. All coefficients
below are smooth on a compact cylinder and $\nu>0$.

\begin{lemma}[Scalar equations]\label{scalar}
The forward equation
\[
 y_t-\nu\Delta y+c\cdot Dy+e y=g,\qquad y(0)=0,
\]
and its time reversal have a unique solution in
$\mathcal W_T=H^1(0,T;L^2)\cap L^2(0,T;H^2)$ for each $g\in L^2$.
For parabolic $C^{\alpha/2,\alpha}$ forcing and $C^{2+\alpha}$
initial or terminal data the solution is in
$C^{1+\alpha/2,2+\alpha}$, with a bounded solution operator.
Writing $B=\|c\|_\infty$ and $E=\|e\|_\infty$, define
\begin{align}
 A&=1+2E+B^2/\nu,&Y&=\sqrt T e^{AT/2},\notag\\
 G&=\sqrt{(1+AY^2)/\nu},&
 J&=\frac{\sqrt3}{\nu}\sqrt{1+B^2G^2+E^2Y^2},\tag{P1}\label{P1}\\
 R(T,\nu,B,E)&=Y+G+J+1+\nu J+BG+EY.\notag
\end{align}
Then $\|y\|_2\leq Y\|g\|_2$, $\|Dy\|_2\leq G\|g\|_2$,
$\|\Delta y\|_2\leq J\|g\|_2$, and
$\|y_t\|_2\leq(1+\nu J+BG+EY)\|g\|_2$.
All these norms are on $Q_T$.
\end{lemma}
\begin{proof}
Multiplication by $y$, Young's inequality, and integration give
\[
 (\|y(t)\|_2^2)' +\nu\|Dy(t)\|_2^2
 \leq \|g(t)\|_2^2+A\|y(t)\|_2^2.
\]
Gronwall gives the bounds $Y,G$. Multiplication by $-\Delta y$ gives
\[
 (\|Dy(t)\|_2^2)' +\nu\|\Delta y(t)\|_2^2
 \leq\frac3\nu(\|g(t)\|_2^2+B^2\|Dy(t)\|_2^2
                                  +E^2\|y(t)\|_2^2).
\]
This proves the bound $J$, and the equation proves the time derivative
bound. Spatial Fourier truncations give finite linear ODEs. These
uniform estimates yield a weak limit with the stated zero trace;
the first estimate applied to a difference proves uniqueness. On a
torus, the $L^2$ norm of the full spatial Hessian equals that of
$\Delta y$, so the asserted space is exactly controlled.

For completeness, the H\"older assertion follows directly from the heat
kernel. If $H_t=e^{\nu t\Delta}$ and
$z(t)=\int_0^t H_{t-s}q(s)\,ds$, the periodic heat kernel satisfies
\[
 \int |D^j H_t(x)|\,\operatorname{dist}(x,0)^\alpha\,dx
       \leq C t^{(\alpha-j)/2}\quad(0<t\leq T),
\]
with the analogous bound for a time derivative (which adds two to $j$).
These follow by differentiating the Gaussian and summing its translates.
The integral of every second spatial derivative is zero. Subtracting
$q(s,x)$ inside the second-derivative convolution therefore gives the
integrable majorant $C\|q\|_\alpha(t-s)^{\alpha/2-1}$.
For a spatial increment $h$, split the integral at $h^2$: the near
part is bounded by $C|h|^\alpha$, and the far part by
$C|h|\int_{h^2}^T r^{\alpha/2-3/2}dr\leq C|h|^\alpha$.
For a time increment $h>0$, splitting at $h$ and using a heat-kernel
time derivative bounds the two parts by
$Ch^{\alpha/2}$ and
$Ch\int_h^T r^{\alpha/2-2}dr\leq Ch^{\alpha/2}$.
Thus $D^2z$ has the required parabolic H\"older bound, and
$z_t=q+\nu\Delta z$ does too. The homogeneous heat solution with
$C^{2+\alpha}$ data satisfies the same estimates by commuting spatial
derivatives and the same kernel bounds.

To include $c\cdot Dy+ey$, use the elementary interpolation estimate
\[
 \|Dy\|_{C^{\alpha/2,\alpha}}+
 \|y\|_{C^{\alpha/2,\alpha}}
 \leq\varepsilon\|y\|_{C^{1+\alpha/2,2+\alpha}}
                         +C_\varepsilon\|y\|_\infty.
\]
It follows by finite spatial differences and splitting increments into
those smaller and larger than a fixed interpolation length; time
increments use the controlled time derivative. Product H\"older
inequalities and a sufficiently small $\varepsilon$ absorb the highest
norm. The maximum principle bounds the remaining supremum norm by
$e^{ET}(\|y(0)\|_\infty+T\|g\|_\infty)$.
Consequently the H\"older estimate is uniform for the coefficients
$sc,se$, $0\leq s\leq1$. The heat equation is solvable at $s=0$.
The bounded inverse estimate makes solvability open by the Neumann
series. It makes solvability closed by compactness in every lower
H\"older exponent and the same estimate for the limiting equation.
This proves existence at $s=1$; the energy estimate proves uniqueness.
Time reversal proves the terminal version. Projecting a scalar
backward solution onto its spatial mean-zero part proves the same
assertion when its equation is spatially centered.
\end{proof}

Classical entire periodic reference solutions in the original statement
are smooth, so the lemma does not strengthen that statement's reference
assumptions. Indeed $m_t,D^2m,v_t,D^2v$ are bounded on the periodic
cylinder. Spatial interpolation shows $Dv$ is $1/2$-H\"older in time
and Lipschitz in space, while $m$ is Lipschitz in time and space.
The HJB right-hand side $f(x,m)-|Dv|^2/2$ is consequently parabolic
$C^{\alpha/2,\alpha}$. Interior versions of the heat-kernel estimates
just proved upgrade $v$. In nondivergence form the population equation
then has H\"older right-hand side $Dv\cdot Dm+m\Delta v$, upgrading
$m$. Differentiating and repeating on interior subcylinders upgrades
both fields to every order. Periodicity covers the entire cylinder
without an initial or terminal regularity exception.

\section{A criterion using scalar problems and finite matrices}
Let $\mathcal H_T=L^2(Q_T)$ with spatial mean zero at almost every time.
For $r\in\mathcal H_T$, first solve the centered scalar terminal problem
and then the scalar initial problem
\begin{equation}\tag{K1}\label{K1}
 \begin{aligned}
 -W_t-\nu\Delta W+\Pi(b\cdot DW)&=\Pi(ar),& W(T)&=0,\\
 Z_t-\nu\Delta Z-\operatorname{div}(bZ)
                         &=\operatorname{div}(\bar m DW),& Z(0)&=0.
 \end{aligned}
\end{equation}
Define $K_{T,\theta}r=Z$. These are two uniquely solvable scalar
equations, rather than an assumed inverse for the coupled MFG system.
Put $B=\|b\|_\infty$, $E=\|\operatorname{div}b\|_\infty$,
$A_0=\|a\|_\infty$, $M=\|\bar m\|_\infty$,
$M_1=\|D\bar m\|_\infty$. Let $G_0,J_0$ be (P1) with $E=0$ and set
\begin{equation}\tag{K2}\label{K2}
 C_K=1+R(T,\nu,B,E)A_0(MJ_0+M_1G_0).
\end{equation}
Lemma~\ref{scalar} proves
\[
 \|Kr\|_2+\|(Kr)_t\|_2+\|\Delta Kr\|_2\leq C_K\|r\|_2.
\]
Thus $K$ is compact: bounded sets in $\mathcal W_T$ are compact in
$L^2$, as is also demonstrated by the explicit projections below.

Let $P_N$ be orthogonal projection onto temporal cosine modes
$0\leq n\leq N$ and spatial Fourier modes
$0<\|k\|_\infty\leq N$. Use real sine--cosine combinations for a real
orthonormal basis $e_1,\ldots,e_{q_N}$ of its range, where
$q_N=(N+1)((2N+1)^d-1)$. Parseval and integration by parts against
the temporal sine modes give
\begin{equation}\tag{K3}\label{K3}
 \|(I-P_N)K\|\leq e_N:=\frac{C_K(T+1)}{N+1}.
\end{equation}
More precisely, the time tail is bounded by
$T\|(Kr)_t\|_2/(\pi(N+1))$ and the spatial tail by
$\|\Delta Kr\|_2/(4\pi^2(N+1)^2)$; their sum implies (K3).
Define the finite real matrix
\begin{equation}\tag{K4}\label{K4}
 D_N=(\delta_{ij}-\langle e_i,Ke_j\rangle)_{i,j=1}^{q_N}.
\end{equation}
Every entry is specified by (K1), using only the original coefficients.
The following is an exact finite-certificate criterion:
\begin{equation}\tag{FC}\label{FC}
 \boxed{\quad\exists N\geq1:\quad \det D_N\ne0,\qquad
 e_N\{1+(1+C_K)\|D_N^{-1}\|_{2\to2}\}<1.\quad}
\end{equation}
No precision or running-time assertion is implicit: arbitrary smooth
coefficients need not even be computably presented. Certified bounds
on the scalar problems and finite matrix entries can verify a strict
certificate when those data are effectively supplied.

\begin{lemma}[Exactness of the finite certificate]\label{finite}
Condition (FC) holds if and only if $I-K$ is invertible on
$\mathcal H_T$.
\end{lemma}
\begin{proof}
In $\operatorname{ran}P_N\oplus\ker P_N$, the operator $I-P_NK$ is
upper triangular, with diagonal blocks $D_N,I$ and upper right block
$-P_NK(I-P_N)$. If $D_N$ is invertible, its full inverse has norm at most
\[
 B_N=1+(1+C_K)\|D_N^{-1}\|.
\]
Since $\|(I-K)-(I-P_NK)\|\leq e_N$, the Neumann series proves
invertibility whenever $e_NB_N<1$.
Conversely, if $I-K$ is invertible, (K3) implies that, for all sufficiently
large $N$, $I-P_NK$ is invertible with inverse norm at most
$2\|(I-K)^{-1}\|$. Restriction to its finite invariant subspace bounds
$\|D_N^{-1}\|$ by the same constant. Thus $B_N$ is bounded and
$e_NB_N<1$ eventually. This proves both directions.
\end{proof}

\section{Fredholm reduction of the Jacobi boundary problem}
The centered linear boundary operator is
\begin{equation}\tag{L1}\label{L1}
 L(w,\rho)=\bigl(\Pi[-w_t-\nu\Delta w+b\cdot Dw-a\rho],\
 \rho_t-\nu\Delta\rho-\operatorname{div}(b\rho+\bar m Dw),\
 \rho(0),w(T)\bigr).
\end{equation}
Its target is the product of two mean-zero parabolic
$C^{\alpha/2,\alpha}$ spaces and $D$. In particular it includes both
interior residuals as well as the two boundary data.

\begin{lemma}[General Fredholm reduction]\label{fredholm}
For every original model, phase and finite horizon, $L$ has Fredholm
index zero. It is invertible if and only if its homogeneous boundary
problem has no nonzero solution, and this is also equivalent to (FC).
\end{lemma}
\begin{proof}
Delete $-a\rho$ from the first component of $L$, obtaining $L_0$.
The scalar lemma solves first the terminal value equation for $w$ and
then the initial population equation for $\rho$; hence $L_0$ is an
isomorphism on the indicated H\"older spaces. The deleted term factors
through the compact embedding
$C^{1+\alpha/2,2+\alpha}\hookrightarrow C^{\alpha/2,\alpha}$.
Thus $L$ is an isomorphism plus a compact operator and has index zero.
The compact-operator alternative gives its invertibility precisely
when its kernel is zero. One elementary justification of that
alternative is to approximate the compact operator in norm by a
finite-rank operator with remainder less than one, invert the remainder
by a Neumann series, and reduce kernel and cokernel to the same finite
square matrix. The compactness here permits finite-rank approximation
by spatial--temporal interpolation, so no Banach approximation
property is presumed.

For a zero-boundary solution, the scalar reduction (K1) says exactly
$\rho=K\rho$, and conversely such a fixed density reconstructs $w$.
This is true in $L^2$ and in the H\"older spaces. These kernels agree.
To see this without a dimension restriction, an $L^2$ kernel first
gives $w,\rho\in\mathcal W_T$. All their spatial derivatives of order
at most two are then in $L^2$. Suppose derivatives through order $j$
are known in $L^2$. Differentiate the backward scalar equation to
order $j$. Its leading forcing is $aD^j\rho$; commutators contain
only already controlled derivatives. The zero terminal trace remains
zero, so the scalar lemma puts $D^jw$ in $\mathcal W_T$.
Next differentiate the forward equation: its principal forcing is
$\operatorname{div}(\bar m D(D^jw))$; commutators contain derivatives
of $\rho$ through order $j$ and of $w$ through order $j+1$, now
controlled. Its zero initial trace again gives
$D^j\rho\in\mathcal W_T$. Spatial difference quotients justify
these distributional differentiations and the zero traces.
This gains two orders and iterates indefinitely. The elementary
trace estimate
$H^1_tH^s_x\cap L^2_tH^{s+2}_x\subset C_tH^{s+1}_x$
(obtained by differentiating the squared $H^{s+1}$ Fourier norm)
then gives continuity of every spatial derivative. The equations
give all time derivatives in succession. Both fields are smooth
up to their respective time endpoints and belong to the H\"older
kernel. Lemma~\ref{finite} completes the equivalence.
\end{proof}

\section{Boundary traces detect every nonzero Jacobi kernel}
\begin{lemma}[Zero at both endpoints]\label{doublezero}
A homogeneous centered Jacobi field $(w,\rho)$ for which
$w(0)=\rho(0)=w(T)=\rho(T)=0$ is identically zero.
\end{lemma}
\begin{proof}
Set $\kappa=\bar m/(2\nu)$, $r=\rho+\kappa w$, $U=(w,r)$ and
$A=\operatorname{diag}(\nu(-\Delta),-\nu(-\Delta))$ on the real
$L^2$ product space. The first component has spatial mean zero; the
second need not. Let
\[
 c_w(t)=\int(b\cdot Dw-a\rho)
       =\int[-(\operatorname{div}b)w-ar+a\kappa w].
\]
Using $\bar m_t=\nu\Delta\bar m+\operatorname{div}(\bar m b)$,
direct substitution cancels all cross second derivatives and gives
\begin{align*}
 w_t&=\nu(-\Delta)w+b\cdot Dw+a\kappa w-ar-c_w,\\
 r_t&=-\nu(-\Delta)r+b\cdot Dr+
           (\operatorname{div}b-a\kappa)r+a\kappa^2w-\kappa c_w.
\end{align*}
Thus $U_t=AU+B(t)U$, where $A$ is a constant self-adjoint operator
and $B(t)$ has only first spatial derivatives and bounded terms.
Integration by parts in the expression for $c_w$ makes its bounded
rank-one character explicit. For some finite $C_0$,
\[
 \|B(t)U\|^2\leq C_0(\langle |A|U,U\rangle+\|U\|^2).
\]
Both endpoint traces of $U$ vanish. Put $V=e^{\phi(t)}U$, where
$\phi'>0$ will be chosen. Integration by parts yields the exact identity
\[
 \int_0^T\|V_t-(A+\phi')V\|^2dt
 =\int_0^T\{\|V_t\|^2+\|(A+\phi')V\|^2
                              +\phi''\|V\|^2\}dt.
\]
The endpoint terms are zero. Spatial Fourier truncation followed by
the $\mathcal W_T$ limit justifies the identity at this regularity.
On the other hand $V_t-(A+\phi')V=BV$, and
\[
 \langle|A|V,V\rangle\leq\|AV\|\|V\|
 \leq\|(A+\phi')V\|\|V\|+\phi'\|V\|^2.
\]
Young's inequality therefore bounds the left side by
\[
 \tfrac12\int_0^T\|(A+\phi')V\|^2dt
            +C_1\int_0^T(1+\phi')\|V\|^2dt.
\]
Choose $\phi'(t)=q e^{kt}$ with $q\geq1$ and $k>2C_1$.
Then $\phi''=k\phi'>C_1(1+\phi')$, so the two displays force
$V=0$. This proves the lemma. The proof uses no forward MFG
semigroup or sign assumption on $a$.
\end{proof}

For a zero-boundary Jacobi field $(z,\sigma)$, namely
$\sigma(0)=0,z(T)=0$, and a perturbation $(w,\rho)$ with interior
linear residuals $(F,G)$, integration by parts gives
\begin{equation}\tag{G1}\label{G1}
 \int d_T\sigma(T)+\int z(0)d_0
       =-\int_{Q_T}(F\sigma+zG).
\end{equation}
Indeed differentiating $\int(w\sigma-z\rho)$ cancels both Laplacians,
both drift terms, and the symmetric coupling terms, leaving
$-\int(F\sigma+zG)$. Mean projections disappear because $z,\sigma$
have mean zero. If both boundary traces $z(0),\sigma(T)$ are zero,
Lemma~\ref{doublezero} says the field is zero. Consequently every
nonzero kernel field gives a nonzero functional on the independent
boundary directions $(d_0,d_T)\in D$.

\section{Exact nonlinear classification and the selection rule}
\begin{theorem}\label{classification}
For every reference allowed in C73-40, every $0<\alpha<1$, every
$T>0$, and every phase $\theta$, the following are equivalent:
\begin{enumerate}
\item Selected local Lipschitz boundary stability (W1) holds.
\item The homogeneous Jacobi problem (L1) has zero kernel.
\item The finite-certificate condition (FC) holds.
\end{enumerate}
When they hold, the rule below selects the unique nearby equilibrium
and is Lipschitz between boundary data; it also satisfies (W2).
When they fail, no rule satisfying (W1) for all sufficiently small
independent boundary perturbations exists, even without continuity
of that rule. A reference has the property on a specified set $H$
exactly when (FC) holds at every pair in $H$.
\end{theorem}
\begin{proof}
The equivalence of items 2 and 3 is Lemma~\ref{fredholm}. Write the
nonlinear perturbation equations as
\[
 L\xi+N(\xi)=(0,0,d_0,d_T),\qquad \xi=(w,\rho),
\]
where
\begin{equation}\tag{N1}\label{N1}
 N(\xi)=\left(\Pi\left[\tfrac12|Dw|^2-
 \{f(x,\bar m+\rho)-f(x,\bar m)-a\rho\}\right],
 -\operatorname{div}(\rho Dw),0,0\right).
\end{equation}
On a sufficiently small positive-density ball in $X_T$, Taylor's
integral formula and the product H\"older inequality give
\begin{equation}\tag{N2}\label{N2}
 \|N(\xi)\|\leq C_2\|\xi\|^2,\qquad
 \|N(\xi)-N(\eta)\|\leq C_2(\|\xi\|+\|\eta\|)\|\xi-\eta\|.
\end{equation}
Explicitly, the scalar Taylor remainder is
$\rho^2\int_0^1(1-s)f_{mm}(x,\bar m+s\rho)ds$; its multiplier
and its Lipschitz dependence have bounded H\"older norms on the fixed
compact positive density interval. The other components are products
of the derivatives controlled by $X_T$.

If item 2 holds, let $K_L=\|L^{-1}\|$. Choose $r>0$ below the
positive-density and declared admissibility margins and with
$2C_2K_Lr\leq1/2$. Take $\delta\leq r/(2K_L)$. Starting at $\xi_0=0$,
perform the explicit iteration
\begin{equation}\tag{SEL}\label{SEL}
 \xi_{n+1}=L^{-1}\{(0,0,d_0,d_T)-N(\xi_n)\}.
\end{equation}
It maps the closed ball of radius $r$ into itself, since its norm
is at most $K_L\|d\|+K_LC_2r^2\leq3r/4$, and its contraction
factor is at most $2K_LC_2r\leq1/2$. Successive differences sum
geometrically, proving convergence and uniqueness in this ball.
Its limit satisfies
\[
 \|\xi(d)\|\leq K_L\|d\|+\tfrac12\|\xi(d)\|,
 \qquad \|\xi(d)-\xi(d')\|\leq2K_L\|d-d'\|.
\]
The first inequality gives (W1), and the second gives the claimed
Lipschitz selection. To implement each linear step, the scalar
reduction solves $I-K$; a certificate (FC) gives the convergent
Neumann series around $I-P_NK$ from Lemma~\ref{finite}. Thus the
rule does not presume an unspecified equilibrium oracle.

The projected first equation has a spatially constant residual
$r_u(t)$ for $u_{\rm r}+w$. Choose a scalar $c(t)$ with
$c'(t)=r_u(t)$ and
$c(T)=\int(h-u_{\rm r}(T))$. Then
$u=u_{\rm r}+w+c$ solves the unprojected HJB equation and has exactly
terminal value $h$. The population equation, initial condition, mass
one, and positivity already hold. This proves actual nonemptiness
for every boundary datum, and (W2) follows by the reference phase
$t+\theta$.

Conversely suppose item 1 holds and a nonzero zero-boundary Jacobi
field $(z,\sigma)$ exists. For any fixed boundary direction
$d\in D$, choose $\varepsilon>0$ small enough that $\varepsilon d$
lies in the open permitted boundary neighborhood. Its selected
perturbation has norm at most $C\varepsilon\|d\|$ and interior
linear residuals $-(N_1,N_2)$. Equation (G1) yields
\[
 \varepsilon\left(\int d_T\sigma(T)+\int z(0)d_0\right)
   =\int_{Q_T}(N_1\sigma+zN_2)=O(\varepsilon^2).
\]
Divide by $\varepsilon$ and let it decrease to zero. Since the
boundary direction was arbitrary, both mean-zero traces
$z(0),\sigma(T)$ vanish. Lemma~\ref{doublezero} contradicts
nonzeroness. This proves necessity, including the negative
classification at every conjugate pair. It rules out exactly (W1),
not every continuous selector or the original orbital property.
\end{proof}

On a compact set of horizon--phase pairs where the criterion holds,
the above constants and the radius may be chosen uniformly. Rescale
time to $[0,1]$; the coefficients and $L$ vary continuously in operator
norm. The inverse Neumann estimate is local in those parameters, so
compactness gives a uniform inverse bound. Reference positivity and
the product estimates are also uniform there. This argument supplies
no uniform bound on an unbounded set of horizons or across a singular
pair. It is compatible with the phase obstruction in the prior attempt.

\section{An entire periodic family: every horizon is classified}
The general criterion has a nonvacuous application to explicit genuine
periodic equilibria. Work on $\mathbb T^1$ and fix
$0<r<1$, $c\ne0$, $\nu>0$. Put $k=2\pi$, $s=\sqrt{1-r^2}$ and
\begin{equation}\tag{TW1}\label{TW1}
 \begin{aligned}
 M(y)&=1+r\cos(ky),\\
 p(y)&=c\left(\frac{s}{M(y)}-1\right)-\nu\frac{M'(y)}{M(y)},\\
 f(z)&=\frac{s^2(c^2+\nu^2k^2)}{2z^2},\qquad
 \lambda=\frac{c^2+\nu^2k^2}{2}.
 \end{aligned}
\end{equation}
Since $\int M^{-1}=s^{-1}$ and $\int M'/M=0$, $p$ has mean zero.
Let $V'=p$ with $\int V=0$. The fields
\begin{equation}\tag{TW2}\label{TW2}
 v_*(t,x)=V(x-ct),\qquad m_*(t,x)=M(x-ct)
\end{equation}
solve the original periodic MFG system, with period $P=1/|c|$.
Indeed, setting $b_0=p+c$ gives
$b_0M=cs-\nu M'$, so $-\nu M''-(b_0M)'=0$.
Furthermore $M''=-k^2(M-1)$ and
$(M')^2=k^2(r^2-(M-1)^2)$ give
\[
 \tfrac12b_0^2-\nu b_0'
 =\frac{s^2(c^2+\nu^2k^2)}{2M^2}-\frac{\nu^2k^2}{2}.
\]
Using $cp+p^2/2=(b_0^2-c^2)/2$ proves
$\lambda+cp-\nu p'+p^2/2=f(M)$, the required HJB equation.
The density is positive, has mass one, and is genuinely time-dependent.
The integral formula for $M^{-1}$ follows, for example, by the
substitution $z=\tan(k y/2)$ on the two half-circles.

\begin{theorem}[All horizons for the explicit family]\label{wave}
For every triple $(r,c,\nu)$ above there is a set
$\Sigma\subset(0,\infty)$, bounded away from zero and finite in every
bounded interval, such that at every phase:
\begin{itemize}
\item For $T\notin\Sigma$, selected local Lipschitz stability (W1)
holds, with the convergent selection (SEL).
\item For $T\in\Sigma$, no selection satisfying (W1) for all small
independent boundary perturbations exists.
\end{itemize}
The set is exactly the zero-kernel failure set in Theorem
\ref{classification}, so (FC) classifies its complement without an
assumed mixed-boundary inverse. The constants can be uniform over
phase and any compact subset of $(0,\infty)\setminus\Sigma$.
\end{theorem}
\begin{proof}
Use $y=x-c(t+\theta)$, $\tau=t/T$, and $q(y)=f'(M(y))$.
The uncentered linear fields satisfy, on the fixed cylinder
$[0,1]\times\mathbb T$,
\begin{equation}\tag{TW3}\label{TW3}
 \begin{aligned}
 B_TW:=-W_\tau-T\nu W_{yy}+Tb_0W_y&=TqR, &W(1)&=0,\\
 F_TR:=R_\tau-T\nu R_{yy}-T(b_0R)_y&=T(MW_y)_y, &R(0)&=0.
 \end{aligned}
\end{equation}
Its density has mean zero. Allowing the scalar value mean here
changes neither its gradient nor the density Schur operator from
(K1); centering the scalar solution recovers exactly that operator.

Let $\mathcal H=L^2([0,1]\times\mathbb T)$,
$\mathcal H_0=\{R:\int R(\tau)=0\}$, and
$\mathfrak D=H^1_\tau L^2_y\cap L^2_\tau H^2_y$.
Use the fixed closed domain $W(1)=0$ for $B_T$, and the fixed
domain $R(0)=0$, $R\in\mathcal H_0$ for $F_T$.
The scalar lemma makes both maps bounded isomorphisms onto their
respective $L^2$ targets for every real $T>0$. Define
\begin{equation}\tag{TW4}\label{TW4}
 \mathcal K_T=T^2F_T^{-1}\partial_y(M\partial_y)B_T^{-1}q
             :\mathcal H_0\longrightarrow\mathcal H_0.
\end{equation}
The derivative operator is bounded from $\mathfrak D$ to $\mathcal H_0$,
and the outer forward inverse maps into $\mathfrak D$, whose embedding
in $\mathcal H$ is compact by the same Fourier--cosine tail estimate
as (K3). Thus $\mathcal K_T$ is compact, including both time endpoints.

On these common domains the scalar operators are affine in $T$.
Around any real $T_0>0$, their inverses are therefore holomorphic
for sufficiently small complex $|T-T_0|$, by the inverse Neumann
series. The union $\Omega$ of such disks is connected: it contains
the entire positive real axis and every disk meets that axis.
The inverse formulas agree on overlaps. Hence $\mathcal K_T$ is
a holomorphic compact family on $\Omega$.

There is an invertible point, rather than an identically singular
family. For $0<T\leq1$, the backward energy estimate for the first
equation of (TW3) gives
\[
 \|W(\tau)\|_2\leq CT\int_\tau^1\|R(s)\|_2ds,
 \quad \|W\|_{L^2}\leq CT\|R\|_{L^2},
 \quad \|W_y\|_{L^2}\leq C\sqrt T\|R\|_{L^2}.
\]
Here and below $C$ depends on the fixed coefficients, not on these
small horizons. The first estimate follows by Gronwall applied to
the spatial $L^2$ norm; testing the equation with $W$ then gives
the gradient estimate. For $Z=\mathcal K_TR$, the forward equation,
tested with $Z$, gives
\[
 (\|Z\|_2^2)' +T\nu\|Z_y\|_2^2
       \leq CT\|Z\|_2^2+CT\|W_y\|_2^2.
\]
Gronwall yields
$\|Z\|_{L^2}\leq C\sqrt T\|W_y\|_{L^2}\leq CT\|R\|_{L^2}$.
Thus $\|\mathcal K_T\|\leq CT<1$ at all sufficiently small $T$.

We include the analytic compact-operator argument. At any
$z_0\in\Omega$, choose a finite-rank orthogonal projection $P$
with $\|(I-P)\mathcal K_{z_0}\|<1/2$. With $Q=I-P$, this stays
less than one near $z_0$, and the $Q$ equation is solved by a
Neumann series. The remaining finite matrix is
\begin{equation}\tag{TW5}\label{TW5}
 I-P\mathcal K_zP-P\mathcal K_zQ
       (I-Q\mathcal K_zQ)^{-1}Q\mathcal K_zP.
\end{equation}
It is holomorphic, and singularity of $I-\mathcal K_z$ is exactly
vanishing of its determinant. In a small connected disk either the
zeros are isolated or the determinant vanishes identically.
The set of points with an identically singular neighborhood is open
and also closed: near a limit point use (TW5); its determinant
vanishes on an open subdisk and hence throughout the disk by the
scalar identity theorem. Connectedness of $\Omega$ and the small-$T$
invertible point show that this open-and-closed set is empty.
All singularities are therefore isolated in $\Omega$.
Their intersection $\Sigma$ with the positive axis is locally finite;
the small-horizon bound keeps it away from zero. The kernel bootstrap
in Lemma~\ref{fredholm} makes every $L^2$ kernel classical.

The moving coordinates give the same problem (TW3) at every phase.
Changing phase changes only a constant spatial translation, an
isometry of the boundary and perturbation norms. Finally Theorem
\ref{classification} gives both the positive and negative nonlinear
conclusions, and its compact-parameter argument proves the final
uniformity assertion.
\end{proof}

\section{What has been classified}
For every original periodic reference, dimension, coupling, phase and
finite horizon, Theorem~\ref{classification} gives a necessary and
sufficient criterion for (W1), including its failure when a Jacobi
kernel is present. The rule (SEL) gives actual equilibria for every
small permitted boundary datum; positivity, mass and the terminal
constant are included. A horizon or phase is not removed from the
classification merely because it is singular. Theorem~\ref{wave}
applies this general result to an explicit genuinely periodic family.

The constants are local to the stated horizon set and the boundary
topology is $C^{2+\alpha}$. No asymptotic attraction, uniform-in-horizon
all-equilibrium stability, finite decision algorithm, computational
complexity bound or formal Lean verification is asserted.

\begin{thebibliography}{9}
\bibitem{CD} P. Cardaliaguet and F. Delarue,
\emph{Selected topics in mean field games}, Proc. ICM 2022, vol.~5,
pp.~3660--3703, EMS Press. DOI:
\href{https://doi.org/10.4171/ICM2022/17}{10.4171/ICM2022/17}.
Section~3.3, ``Periodic solutions,'' printed p.~3681.
\url{https://ems.press/content/book-chapter-files/33265}.
\bibitem{CN} M. Cirant and L. Nurbekyan,
\emph{The variational structure and time-periodic solutions for
mean-field games systems}, Minimax Theory Appl.~3(2) (2018), 227--260.
Theorem~6.1, arXiv v1 pp.~24--25; Remark~6.4, p.~28.
\url{https://arxiv.org/pdf/1804.08943v1}.
\end{thebibliography}
\end{document}
